% !TEX root = ../main.tex
% !TEX program = lualatex
\documentclass[../main.tex]{subfiles}
\IfSubfilesClassLoaded{\externaldocument{\subfix{../build/main}}}{}
% =============================================================================
% File: 23_Waterfront_Organization_Work_Control_Safety.tex
% =============================================================================
\begin{document}
\IfSubfilesClassLoaded{\chapter{EDO Study Guide}}{}
\section{Waterfront Organizations, Work Controls, and Safety}

\subsection{Summary}
\begin{description}
  \item[\textbf{Waterfront focus.}] Most \ac{syscom} products are installed during new construction or overhaul periods, so program managers must understand waterfront safety and work control rules~\autocite{edo-4-1-5-waterfront-work-control-2025}.
  \item[\textbf{Key nuclear players.}] \ac{navreact}, \ac{nrro}, \ac{rpco}, and Joint Test Groups enforce nuclear safety and testing discipline at public shipyards~\autocite{edo-4-1-5-waterfront-work-control-2025}.
  \item[\textbf{Safety programs.}] The loss of \ac{uss} THRESHER drove \ac{subsafe} requirements, and the loss of \ac{uss} GRAYBACK established the \ac{dss} program~\autocite{edo-4-1-5-waterfront-work-control-2025}.
  \item[\textbf{Industrial safety manuals.}] The GUITARRO accident drove the 6010 manual for submarine ship safety, while the MIAMI fire drove the 8010 manual for fire prevention and response~\autocite{edo-4-1-5-waterfront-work-control-2025}.
  \item[\textbf{Work control.}] All work requires a valid \ac{twd} and authorization, with tag-out requirements defined by the \ac{tum}~\autocite{edo-4-1-5-waterfront-work-control-2025}.
\end{description}

\subsection{Waterfront Context}
\begin{itemize}
  \item \ac{syscom} products are installed at the waterfront during new construction, overhaul, and repair periods~\autocite{edo-4-1-5-waterfront-work-control-2025}.
  \item Waterfront culture emphasizes safety of people, equipment, and facilities; program managers may be held accountable for compliance~\autocite{edo-4-1-5-waterfront-work-control-2025}.
\end{itemize}

\subsection{Naval Nuclear Propulsion Program}
\begin{description}
  \item[\textbf{Program basis.}] The Atomic Energy Act of 1954 established the Naval Nuclear Propulsion Program as a joint Department of Energy and Navy effort, with Executive Order 12344 and Public Law 98-525 defining responsibilities, and the FY2000 NDAA relocating Naval Reactors under the National Nuclear Security Administration~\autocite{edo-4-1-5-waterfront-work-control-2025}.
  \item[\textbf{\ac{nrro}.}] Ensures reactor safeguards are not compromised by enforcing standards, providing independent assessment, and keeping Naval Reactors informed~\autocite{edo-4-1-5-waterfront-work-control-2025}.
  \item[\textbf{\ac{rpco}.}] Represents prime reactor contractors, supports joint test and refueling groups, and can stop work or testing if unsafe conditions exist~\autocite{edo-4-1-5-waterfront-work-control-2025}.
  \item[\textbf{Joint Test Groups.}] Coordinate nuclear and non-nuclear testing with the \ac{nsa} chief test engineer, \ac{nrro} oversight, Ship's Force participation, and reactor contractor technical advice~\autocite{edo-4-1-5-waterfront-work-control-2025}.
\end{description}

\subsection{SUBSAFE and Deep Submergence Systems}
\begin{description}
  \item[\textbf{Origins.}] The loss of \ac{uss} THRESHER led to the \ac{subsafe} program, and the loss of \ac{uss} GRAYBACK led to the \ac{dss} program~\autocite{edo-4-1-5-waterfront-work-control-2025}.
  \item[\textbf{SUBSAFE purpose.}] The \ac{subsafe} program provides maximum reasonable assurance of hull integrity and operability of critical systems needed to control and recover from flooding casualties~\autocite{edo-4-1-5-waterfront-work-control-2025}.
  \item[\textbf{Fundamentals.}] Work discipline, material control, documentation (design products and \ac{oqe}), and compliance verification through inspections, reviews, audits, and certification~\autocite{edo-4-1-5-waterfront-work-control-2025}.
  \item[\textbf{General requirements.}] \ac{subsafe} certification is maintained through life via re-entry control, the \ac{uro} program, and certification and functional audits; initial certification is coordinated through the cognizant \ac{supship}~\autocite{edo-4-1-5-waterfront-work-control-2025}.
  \item[\textbf{Level I material.}] Level I material requires testing, traceability to original test reports, and receipt inspection by \ac{navsea}-approved certifying activities~\autocite{edo-4-1-5-waterfront-work-control-2025}.
  \item[\textbf{\ac{dss} scope.}] Covers manned noncombatant submersibles, submarine-oriented \ac{dss} including diving systems, and man-rated handling systems, applicable when Navy personnel use any \ac{dss} or any personnel use a Navy-owned \ac{dss}~\autocite{edo-4-1-5-waterfront-work-control-2025}.
  \item[\textbf{Fly-by-wire program.}] Provides maximum reasonable assurance that submarine fly-by-wire controls will not cause or impede recovery from a jam dive casualty; requires re-entry controls within the flight critical boundary with \ac{isea} support~\autocite{edo-4-1-5-waterfront-work-control-2025}.
\end{description}

\subsection{Industrial Ship Safety Manuals}
\begin{description}
  \item[\textbf{6010 manual (GUITARRO).}] The \ac{uss} GUITARRO sinking led to the 6010 manual, which controls work and testing that affect ship conditions during construction and availabilities, especially for buoyancy, list, trim, watertight integrity, and stability~\autocite{edo-4-1-5-waterfront-work-control-2025}.
  \item[\textbf{6010 requirements.}] Establishes ship safety and ship test organizations, controlling documents (safety plan of the day, daily test schedule, and prerequisite lists), and flooding and fire prevention procedures~\autocite{edo-4-1-5-waterfront-work-control-2025}.
  \item[\textbf{8010 manual (MIAMI).}] The \ac{uss} MIAMI fire led to the 8010 manual for fire prevention and response in the industrial environment and highlighted the need for early detection, rapid response, and extended response capability~\autocite{edo-4-1-5-waterfront-work-control-2025}.
  \item[\textbf{8010 applicability.}] Applies to all ship repair and construction activities and ship availabilities, requiring \ac{moa} agreements on fire prevention, detection, and response roles~\autocite{edo-4-1-5-waterfront-work-control-2025}.
  \item[\textbf{Fire safety governance.}] Fire safety councils coordinate local approvals for hull cuts, fire zones, access routes, alarms, and firefighting systems, and oversee compliance and risk mitigations~\autocite{edo-4-1-5-waterfront-work-control-2025}.
\end{description}

\subsection{Work Authorization and Control}
\begin{description}
  \item[\textbf{Work control triangle.}] Training, procedures, supervision, and audit or \ac{oqe} scale with the complexity and criticality of the work~\autocite{edo-4-1-5-waterfront-work-control-2025}.
  \item[\textbf{Unplanned events.}] Team learning sessions (submarines) or critiques (surface) drive root cause analysis, corrective actions, and adjustment of controls when severity is unacceptable~\autocite{edo-4-1-5-waterfront-work-control-2025}.
  \item[\textbf{Minimum requirements.}] All work requires a valid \ac{twd} and authorization to perform work~\autocite{edo-4-1-5-waterfront-work-control-2025}.
  \item[\textbf{\ac{twd} types.}] Maintenance procedures, formal work packages, and controlled work packages are selected based on task complexity, worker knowledge, and the need for \ac{oqe}~\autocite{edo-4-1-5-waterfront-work-control-2025}.
  \item[\textbf{Authorization process.}] All outside activity work requires formal authorization via the \ac{waf} to ensure safety standards in public and private availabilities~\autocite{edo-4-1-5-waterfront-work-control-2025}.
\end{description}

\subsection{Tag-Outs in Maintenance}
\begin{description}
  \item[\textbf{TUM guidance.}] The \ac{tum} covers repair activity responsibilities, barrier criteria, divers' tags, and maintenance-period tag-out situations, with specific appendices for amplifications and shift operations~\autocite{edo-4-1-5-waterfront-work-control-2025}.
  \item[\textbf{New construction.}] Contractors use occupational safety and health isolation processes until the pre-commissioning crew embarks or, for nuclear work, when the crew takes control of a system~\autocite{edo-4-1-5-waterfront-work-control-2025}.
  \item[\textbf{Repair activity responsibilities.}] Ensure qualified personnel, verify tag-out accuracy, sign tag-outs as the repair activity representative, and remove tags when no longer needed~\autocite{edo-4-1-5-waterfront-work-control-2025}.
  \item[\textbf{Divers' tags.}] Diving activity representatives review tags, brief supervisors on isolations, and coordinate with \acp{waf} when multiple tag-out logs are affected~\autocite{edo-4-1-5-waterfront-work-control-2025}.
  \item[\textbf{Double barrier protection.}] Required for high temperature, high pressure, sea-connected systems, hull penetrations below the waterline, low flash point fluids, oxygen, and hazardous or toxic vapors~\autocite{edo-4-1-5-waterfront-work-control-2025}.
\end{description}

\subsection{Quick Review}
\begin{itemize}
  \item Naval Reactors' primary mission is reactor safety; \ac{rpco} reports to Department of Energy Naval Reactors~\autocite{edo-4-1-5-waterfront-work-control-2025}.
  \item THRESHER led to \ac{subsafe} requirements, while \ac{dss} protects occupants of deep submergence systems~\autocite{edo-4-1-5-waterfront-work-control-2025}.
  \item GUITARRO drove the 6010 manual; MIAMI drove the 8010 manual for industrial fire safety~\autocite{edo-4-1-5-waterfront-work-control-2025}.
  \item Minimum work requirements are a valid \ac{twd} and authorization; double-barrier conditions follow \ac{tum} guidance~\autocite{edo-4-1-5-waterfront-work-control-2025}.
\end{itemize}

\IfSubfilesClassLoaded{
  \RenewDocumentCommand{\entryname}{}{\textbf{\color{Modern} Acronym}}
  \RenewDocumentCommand{\descriptionname}{}{\textbf{\color{Modern} Definition}}
  \printnoidxglossary[
    type=\acronymtype,
    title=Acronyms,
    style=long-booktabs
  ]}{}
\subsection{Coursebook cue: RPCO and Deep Submergence Systems}

The RPCO is a Naval Nuclear Laboratory waterfront presence that oversees
technical aspects of naval nuclear work and evolutions.  The Resident Manager,
RPCO has authority to stop work, operations, or testing when conditions are
unsafe, and the RPCO can provide technical advice to nuclear test groups even
when not acting as a voting member
\autocite{edo-4-1-5-waterfront-work-control-2025}.

The Deep Submergence Systems program grew from the need to protect occupants of
deep-submergence systems from material or procedural failures.  It applies to
Navy-owned or Navy-used deep-submergence systems, including manned
noncombatant submersibles, submarine-oriented diving systems, and man-rated
portions of handling systems.  The board cue is that DSS is a certification and
safety-assurance regime, not a generic diving label
\autocite{edo-4-1-5-waterfront-work-control-2025}.

\end{document}
